%% Response to Editor — INS-D-26-5725
%% Fisher--Rao Geometric Optimization and Detection for
%% Generative Steganographic Transport
%% Ekodeck et al. — Information Sciences

\documentclass[12pt,a4paper]{article}

\usepackage[utf8]{inputenc}
\usepackage[T1]{fontenc}
\usepackage[margin=2.5cm]{geometry}
\usepackage{amsmath,amssymb}
\usepackage{booktabs}
\usepackage{tabularx}
\usepackage{array}
\usepackage{xcolor}
\usepackage{enumitem}
\usepackage{hyperref}
\usepackage{microtype}
\usepackage{parskip}
\usepackage{titlesec}
\usepackage{fancyhdr}
\usepackage{mdframed}

% ── Colours ──────────────────────────────────────────────────────────────────
\definecolor{editorblue}{RGB}{0,60,120}
\definecolor{responsegreen}{RGB}{0,90,50}
\definecolor{shadeblue}{RGB}{235,242,252}
\definecolor{shadegreen}{RGB}{235,248,240}

% ── Section formatting ────────────────────────────────────────────────────────
\titleformat{\section}[block]
  {\large\bfseries\color{editorblue}}{\thesection.}{0.5em}{}[\vspace{2pt}\hrule\vspace{4pt}]
\titleformat{\subsection}[block]
  {\normalsize\bfseries\color{editorblue}}{\thesubsection}{0.5em}{}

% ── Header / footer ───────────────────────────────────────────────────────────
\pagestyle{fancy}
\fancyhf{}
\rhead{\small\color{gray} INS-D-26-5725 — Response to Editor}
\lhead{\small\color{gray} Ekodeck et al.}
\rfoot{\small\thepage}
\renewcommand{\headrulewidth}{0.4pt}

% ── Environments ──────────────────────────────────────────────────────────────
\newmdenv[
  backgroundcolor=shadeblue,
  linecolor=editorblue,
  linewidth=1.2pt,
  leftline=true, rightline=false, topline=false, bottomline=false,
  innerleftmargin=10pt, innerrightmargin=8pt,
  innertopmargin=6pt, innerbottommargin=6pt,
  skipabove=8pt, skipbelow=6pt
]{editorbox}

\newmdenv[
  backgroundcolor=shadegreen,
  linecolor=responsegreen,
  linewidth=1.2pt,
  leftline=true, rightline=false, topline=false, bottomline=false,
  innerleftmargin=10pt, innerrightmargin=8pt,
  innertopmargin=6pt, innerbottommargin=6pt,
  skipabove=4pt, skipbelow=6pt
]{responsebox}

% ── Macros ────────────────────────────────────────────────────────────────────
\newcommand{\DKL}{D_{\mathrm{KL}}}
\newcommand{\dg}{d_g}

% ─────────────────────────────────────────────────────────────────────────────

\begin{document}

% ── Title block ───────────────────────────────────────────────────────────────
\begin{center}
  {\LARGE\bfseries\color{editorblue} Response to the Editor}\\[8pt]
  {\large\bfseries Information Sciences}\\[4pt]
  {\large Manuscript No.\ \textbf{INS-D-26-5725}}\\[6pt]
  \rule{\linewidth}{0.8pt}\\[4pt]
  {\large\itshape Fisher--Rao Geometric Optimization and Detection\\
  for Generative Steganographic Transport}\\[6pt]
  {\normalsize
    St\'ephane Ga\"el R.\ Ekodeck,
    Prosper Rosaire Mama Assandje,
    Serge Alain \'Eb\'el\'e,\\
    St\'ephane Willy Mossebo Tcheunteu,
    Arthur Ulrich Ewane,
    Damaris Belle M.\ Fotso,
    Ren\'e Ndoundam
  }\\[4pt]
  {\small Submitted: 2026 \quad|\quad Revised: \today}
  \rule{\linewidth}{0.8pt}
\end{center}

\vspace{6pt}

Dear Editor,

We are grateful for the careful engagement with our manuscript and for
the opportunity to resubmit a revised version.
Below we address every point raised in full, describe precisely where the
manuscript has been modified, and include a summary table of all changes.

% ─────────────────────────────────────────────────────────────────────────────
\section{Editorial Concern and Response}
% ─────────────────────────────────────────────────────────────────────────────

\subsection{Core editorial concern}

\begin{editorbox}
\textbf{Editor's question (condensed):}
``The use of Fisher--Rao geometry, natural-gradient methods, and
information-geometric optimization frameworks are themselves well-established
mathematical tools.
What is it that makes it different from the previous work?
The manuscript did not demonstrate a sufficiently compelling advancement
relative to the current state of the literature and the high novelty
threshold maintained by \textit{Information Sciences}.''
\end{editorbox}

\begin{responsebox}
\textbf{Our response.}
We fully accept this as the right question and have restructured the
introduction to answer it directly and explicitly.
The tools (Fisher--Rao metric, $\alpha$-connections, natural gradient) are
indeed established.
Our contribution is not that these tools exist, but that, when applied to the
\emph{steganographic transport problem}, they yield concrete results that are
neither obvious nor previously known.
We detail four such results below.
\end{responsebox}

% ─────────────────────────────────────────────────────────────────────────────
\section{What the Framework Produces That Is Genuinely New}
% ─────────────────────────────────────────────────────────────────────────────

\subsection{A new problem formulation}

\begin{editorbox}
\textbf{Gap in the prior literature.}
Every steganographic method prior to this work --- from Cachin's
$\varepsilon$-security criterion to state-of-the-art coverless schemes
iMEC and Discop --- treats embedding as a \emph{point-to-point} problem: it
minimises the endpoint divergence $\DKL(P_{\theta_1}\|P_{\theta_0})$ between
the final stego distribution and the cover distribution, while ignoring how
the generator parameter evolves from $\theta_0$ to $\theta_1$.
\end{editorbox}

\begin{responsebox}
We are the first to ask: \emph{does the trajectory through parameter space
matter for the accumulated detectability of a multi-step embedding process?}
This question has no precursor in either the steganography or the
information-geometry literature.
The answer, formalised by the sequential-KL path-cost functional
\[
  \mathcal{C}[\gamma]
  \;=\;
  \sum_{k=0}^{N-1}
  \DKL\!\bigl(P_{\gamma(t_{k+1})}\bigm\|P_{\gamma(t_k)}\bigr),
\]
is yes: the accumulated detectability depends on the entire trajectory, and
the optimal path is \emph{not} the Euclidean straight line.
The formulation of this path-cost problem is a conceptual contribution
independent of the tools used to solve it.
\end{responsebox}

\subsection{Geodesic optimality (Theorem~2, part~iv): a non-trivial result}

\begin{editorbox}
\textbf{Possible confusion:}
``Geodesics minimise length'' is a tautology. Is Theorem~2 anything more?
\end{editorbox}

\begin{responsebox}
Yes.
The steganographic objective --- minimising $\mathcal{C}[\gamma]$ over a
discretised path --- is \emph{not} the Riemannian length functional.
These two quantities coincide only in the limit of infinitesimally fine
discretisation and under a specific second-order KL expansion.

Theorem~2(iv) proves the following chain:
\begin{enumerate}[nosep]
  \item The local KL expansion gives
        $\DKL(P_{\theta+\delta}\|P_{\theta})
         = \tfrac{1}{2}\delta^\top F(\theta)\,\delta + O(\|\delta\|^3)$.
  \item Summing over a discretised path, the leading-order term is the
        Fisher path-energy
        $E[\gamma]=\int_0^1\dot\gamma^\top F(\gamma)\dot\gamma\,\mathrm{d}t$.
  \item Minimising $E[\gamma]$ over paths with fixed endpoints yields the
        geodesic equation as the Euler--Lagrange condition.
\end{enumerate}
This chain links the steganographic path cost to the Riemannian structure
through a specific asymptotic argument.
It does not appear in the natural-gradient or information-geometry literature
(where the path-cost problem for KL accumulation has not been studied) nor
in the steganography literature (which does not consider trajectory geometry).
A \textbf{Remark} to this effect has been added directly before Theorem~2 in
the revised manuscript.
\end{responsebox}

\subsection{Detection Duality (Theorem~7): a structural result specific to steganography}

\begin{editorbox}
\textbf{Possible confusion:}
The dual-affine structure of exponential families ($e$-geodesics, $m$-projections)
is classical (Amari--Nagaoka). Is Theorem~7 anything more?
\end{editorbox}

\begin{responsebox}
Yes.
The classical result is that exponential families admit a dually flat structure.
What is new is the \emph{adversarial assignment} of the two geometries:
\begin{itemize}[nosep]
  \item The \textbf{embedder}'s optimal path lives in the
        \emph{Levi-Civita} ($\nabla^{(0)}$, Fisher--Rao) geometry ---
        governed by Theorem~2.
  \item The \textbf{optimal Neyman--Pearson steganalyzer}'s rejection region
        is affine in expectation coordinates, i.e.\ it is an \emph{$m$-flat}
        half-space governed by the $m$-geometry ($\nabla^{(-1)}$).  Its normal
        direction is the $e$-geodesic connecting cover and stego parameters
        --- governed by Theorem~7.
\end{itemize}
The Amari--Chentsov $\pm 1$ connection pair thus plays \emph{opposite and
complementary roles for the two adversaries}.
This structural $\pm 1$ duality is specific to the steganographer--steganalyzer
adversarial setting and has no counterpart in prior work on exponential
families or on steganographic security.
An introductory paragraph making this explicit has been added before
Theorem~7 in the revised manuscript.
\end{responsebox}

\subsection{Stein exponent (Theorem~8): closing the loop between embedding and detection geometry}

\begin{editorbox}
\textbf{Possible confusion:}
Stein's lemma identifying $-\DKL(P_{\theta_s}\|P_{\theta_c})$ as the
optimal type-II error exponent is classical.
\end{editorbox}

\begin{responsebox}
The classical Stein result is well known.
What Theorem~8 adds in the steganographic context is the \emph{chain of
connections}:
\[
  \beta_n^{(\alpha)}
  \;\lesssim\;
  \exp\!\Bigl(-\tfrac{n}{2}\,\dg(\theta_c,\theta_s)^2\Bigr),
\]
obtained by expressing the KL divergence through the local Fisher--Rao
distance expansion (Theorem~6).
This closes the loop between the two adversaries:
\begin{itemize}[nosep]
  \item A small Fisher--Rao distance $\dg(\theta_c,\theta_s)$ buys
        low accumulated embedding distortion (Theorem~2) \emph{and}
        high type-II error for the steganalyzer (Theorem~8).
  \item The same geometric quantity governs both the embedding
        cost and the detection difficulty.
\end{itemize}
This unification of embedding geometry and detection geometry through a single
Fisher--Rao distance is new.
It does not appear in the endpoint-only steganographic security literature
(Cachin, Fridrich--Ker), which provides no connection to the trajectory.
A \textbf{Remark} to this effect has been added before Theorem~8 in the
revised manuscript.
\end{responsebox}

\subsection{The two faces of curvature: a new adversarial trade-off}

\begin{responsebox}
In the regime of negative sectional curvature (which includes the Gaussian
family, $\kappa=-1/2$), curvature plays \emph{opposite roles for the two
adversaries simultaneously}:
\begin{itemize}[nosep]
  \item \textbf{Detection advantage:} negative curvature expands the
        Fisher-ball volume accessible to the steganalyzer (amplified
        Jacobi-field growth), increasing the resolution of the detection test.
  \item \textbf{Embedding cost:} the same negative curvature amplifies
        geodesic transport instability for the embedder.
\end{itemize}
This geometric trade-off, intrinsic to the adversarial game, is a new
observation with no precursor in either the steganography or the Riemannian
geometry literature.
\end{responsebox}

% ─────────────────────────────────────────────────────────────────────────────
\section{Summary of All Changes to the Manuscript}
% ─────────────────────────────────────────────────────────────────────────────

\begin{table}[h!]
\centering
\caption{Complete list of modifications in the revised manuscript.}
\label{tab:changes}
\small
\renewcommand{\arraystretch}{1.35}
\begin{tabularx}{\linewidth}{p{3.2cm} X p{4.2cm}}
\toprule
\textbf{Location} & \textbf{Nature of change} & \textbf{Addresses} \\
\midrule
\S1.4 (new subsection)
  & ``Beyond the Tools: What the Framework Produces'' ---
    4 focused paragraphs answering the editor's central question.
  & Core editorial concern \\
Table~1 (new)
  & 10-row $\times$ 5-column positioning table: contributions vs.\ prior
    literature (Cachin/Fridrich--Ker; iMEC/Discop/SparSamp; Amari IG;
    natural-gradient optimisation; this paper).
    Four exclusive contributions are in bold.
  & Core editorial concern \\
Before Theorem~2
  & \texttt{Remark} explaining why part~(iv) is not a tautology of
    ``geodesics minimise length,'' with explicit proof sketch of the
    second-order KL $\to$ Fisher-energy chain.
  & Non-triviality of key results \\
Before Theorem~7
  & Introductory paragraph clarifying that the novelty is not the classical
    dual-affine structure but its adversarial steganographic interpretation
    ($\pm1$ geometry assigned to embedder vs.\ steganalyzer).
  & Non-triviality of key results \\
Before Theorem~8
  & \texttt{Remark} distinguishing from classical Stein's lemma and
    explaining the new chain: Stein exponent $\leftrightarrow$ Fisher--Rao
    distance $\leftrightarrow$ embedding cost.
  & Non-triviality of key results \\
\bottomrule
\end{tabularx}
\end{table}

\noindent
\textbf{No theorems, proofs, simulations, or conclusions have been altered.}
The scope statement (theoretical foundation, not operational pipeline) is
unchanged.

% ─────────────────────────────────────────────────────────────────────────────
\section{Positioning vs.\ Prior Literature}
% ─────────────────────────────────────────────────────────────────────────────

To make the novelty self-evident, the revised manuscript now includes
Table~1 (reproduced below as Table~\ref{tab:novelty_resp}).
Rows in \textbf{bold} represent capabilities exclusive to this paper.

\begin{table}[h!]
\centering
\caption{Positioning of contributions relative to prior literature
  (\checkmark~=~provided; $\circ$~=~partially; $\times$~=~not provided).}
\label{tab:novelty_resp}
\scriptsize
\renewcommand{\arraystretch}{1.35}
\begin{tabularx}{0.98\linewidth}{p{3.6cm} *{5}{>{\centering\arraybackslash}X}}
\toprule
\textbf{Capability}
  & \textbf{Cachin / F-K}
  & \textbf{iMEC / Discop}
  & \textbf{Amari IG}
  & \textbf{Nat.\ grad.}
  & \textbf{This paper} \\
\midrule
KL endpoint security criterion       & \checkmark & \checkmark & $\times$ & $\times$ & \checkmark \\
Coverless / generative embedding     & $\times$ & \checkmark & $\times$ & $\times$ & \checkmark \\
Fisher--Rao metric on param.\ space  & $\times$ & $\times$ & \checkmark & \checkmark & \checkmark \\
Natural-gradient optimisation        & $\times$ & $\times$ & $\circ$  & \checkmark & \checkmark \\
\midrule
\textbf{Path-cost formulation}       & $\times$ & $\times$ & $\times$ & $\times$ & \checkmark \\
\textbf{Sequential-KL geodesic optimality (Thm.~2)} & $\times$ & $\times$ & $\times$ & $\times$ & \checkmark \\
\textbf{Detection-duality ($\pm1$ structure, Thm.~7)} & $\times$ & $\times$ & $\times$ & $\times$ & \checkmark \\
\textbf{Curvature trade-off in adversarial game} & $\times$ & $\times$ & $\times$ & $\times$ & \checkmark \\
\midrule
Convergence guarantees for stego optimizer & $\times$ & $\times$ & $\times$ & \checkmark & \checkmark \\
Numerical illustrations on stat.\ manifolds & $\times$ & $\circ$ & $\times$ & $\circ$ & \checkmark \\
\bottomrule
\end{tabularx}
\end{table}

% ─────────────────────────────────────────────────────────────────────────────
\section*{Closing Remarks}
% ─────────────────────────────────────────────────────────────────────────────

We believe the revised manuscript now answers the editorial question directly,
explicitly, and with concrete precision at every relevant location in the text.
The four contributions exclusive to this paper (bold rows in
Table~\ref{tab:novelty_resp}) are non-trivial consequences of applying
established tools to a genuinely new problem formulation --- the
steganographic transport path problem --- and to a genuinely new adversarial
setting in which the $\pm1$ dual geometry plays structurally distinct roles for
the two adversaries.

We hope the revised manuscript meets the novelty threshold expected by
\textit{Information Sciences} and look forward to your assessment.

\bigskip
Respectfully yours,

\medskip
\textbf{St\'ephane Ga\"el R.\ Ekodeck}\\
Corresponding author\\
University of Yaound\'e~I --- LIRIMA/GRIMCAPE \& UMMISCO, Sorbonne Universit\'es\\
\href{mailto:stephane-gael.ekodeck@facsciences-uy1.cm}
     {stephane-gael.ekodeck@facsciences-uy1.cm}

\medskip
\textit{On behalf of all co-authors:}\\
P.R.\ Mama Assandje,
S.A.\ \'Eb\'el\'e,
S.W.\ Mossebo Tcheunteu,
A.U.\ Ewane,
D.B.M.\ Fotso,
R.\ Ndoundam.

\end{document}
